problem:  
Let \(f: (M^m, g) \to (N^n, h)\) be a **submersive harmonic morphism** between compact oriented Riemannian manifolds, with connected fibers and horizontal distribution \(\mathcal{H}\).  
By standard theory (Fuglede–Ishihara, Baird–Wood), such maps are horizontally conformal and have minimal fibers, with curvature decomposed via the O’Neill tensors \(A\) and \(T\).  

**Rigidity Problem:**  
Assume that \(M\) has **positive Ricci curvature**.  
Is it true that under these assumptions the **O’Neill tensor \(A\)** of \(f\) must vanish identically, i.e. that the horizontal distribution \(\mathcal{H}\) is **integrable**, forcing \(M\) to be (locally) a Riemannian product submersion with totally geodesic minimal fibers?  

Equivalently:  
Does positive Ricci curvature forbid the existence of **nontrivial harmonic morphisms** (those with nonvanishing \(A\)-tensor) between compact Riemannian manifolds?  

Why is it a "good" problem:  
This revised problem is mathematically sound and conceptually simple: it asks whether positive Ricci curvature forces complete rigidity of the horizontal distribution in a harmonic morphism. It removes redundant hypotheses and precisely targets the fundamental, still open, question of interplay between curvature and the submersion’s analytic structure.  

It is “good” because it unites multiple deep themes:  

1. **Analysis:** Harmonic morphisms are nonlinear elliptic systems satisfying overdetermined PDEs expressing conformality and energy-minimizing conditions.  
2. **Geometry:** The O’Neill tensors \(A\) and \(T\) measure how far a submersion deviates from being a Riemannian product. Their vanishing is equivalent to integrability or total geodesicity of fibers.  
3. **Global structure:** Positive Ricci curvature imposes strong topological constraints (e.g., via Myers’ theorem). If such curvature also forbids nontrivial twisting (nonzero \(A\)), that would uncover a new rigidity paradigm coupling analytic, geometric, and topological properties.  

The problem’s appeal lies in its **extreme simplicity of statement** yet **high potential depth:** it asks whether Ricci positivity — a purely curvature condition — annihilates one of the most geometric obstructions, the \(A\)-tensor, for harmonic morphisms. A resolution would likely require new techniques blending Bochner-type identities with the Gray–O’Neill submersion formulas, providing a profound bridge between harmonic analysis and global Riemannian geometry.